
\documentclass{article}
\usepackage{colortbl}
\usepackage{makecell}
\usepackage{multirow}
\usepackage{supertabular}

\begin{document}

\newcounter{utterance}

\centering \large Interaction Transcript for game `ifeval', experiment `ifeval', episode 54 with llama{-}3.1{-}8b{-}instruct.
\vspace{24pt}

{ \footnotesize  \setcounter{utterance}{1}
\setlength{\tabcolsep}{0pt}
\begin{supertabular}{c@{$\;$}|p{.15\linewidth}@{}p{.15\linewidth}p{.15\linewidth}p{.15\linewidth}p{.15\linewidth}p{.15\linewidth}}
    \# & \multicolumn{2}{c}{Player} && \multicolumn{2}{c}{Game Master} \\
    \hline

    \theutterance \stepcounter{utterance}  
    & & & \multicolumn{4}{p{0.6\linewidth}}{
        \cellcolor[rgb]{0.9,0.9,0.9}{
            \makecell[{{p{\linewidth}}}]{
                \texttt{\tiny{[P1$\langle$GM]}}
                \texttt{Take the text below as a starting point, and make it a complete article: "You may have to meet with a helper to work out a parenting plan. The first would be to estimate how many time you have everyday for parenting, and is that enough...."} \\
\texttt{Avoid using the following keywords: sleep, cook, feed} \\
\texttt{Mention the keyword 'schedule' for more than 5 times.} \\
            }
        }
    }
     \\ \\

    \theutterance \stepcounter{utterance}  
    & \multicolumn{4}{p{0.6\linewidth}}{
        \cellcolor[rgb]{0.561,0.737,0.561}{
            \makecell[{{p{\linewidth}}}]{
                \texttt{\tiny{[P1$\rangle$GM]}}
                \texttt{Creating a Parenting Plan: Understanding Your Daily Schedule} \\
\\ 
\texttt{When parents decide to separate or divorce, one of the most critical aspects to consider is creating a parenting plan. This plan outlines how you and your co{-}parent will share responsibilities and make decisions about your child's care. Meeting with a helper, such as a mediator or counselor, can be beneficial in working out a plan that suits your family's unique needs.} \\
\\ 
\texttt{The first step in creating a parenting plan is to estimate how many hours you have available each day for parenting. This involves taking a close look at your daily schedule and determining if it's sufficient to meet your child's needs. Consider your work schedule, personal appointments, and other commitments to get a clear picture of your available time.} \\
\\ 
\texttt{To create a realistic schedule, it's essential to be honest about your daily routine. Ask yourself questions like: What time do I wake up and go to bed? How many hours do I spend at work or school? Are there any regular appointments or activities that I need to attend? By understanding your daily schedule, you can determine how much time you have available for parenting and whether it's enough to meet your child's needs.} \\
\\ 
\texttt{Once you have a clear understanding of your daily schedule, you can start thinking about how to allocate your time effectively. This may involve creating a schedule that outlines specific times for parenting responsibilities, such as spending time with your child, attending school events, or handling household chores. A well{-}planned schedule can help ensure that your child receives the care and attention they need, even when you're not physically present.} \\
\\ 
\texttt{When creating a parenting plan, it's also essential to consider your child's schedule. This may involve coordinating with your co{-}parent to ensure that your child's needs are met, even when you're not together. For example, you may need to create a schedule that outlines who will drop off and pick up your child from school, or who will attend parent{-}teacher conferences.} \\
\\ 
\texttt{A parenting plan should be flexible and adaptable to changing circumstances. Life is unpredictable, and schedules can change unexpectedly. By creating a schedule that is realistic and flexible, you can ensure that your child's needs are met, even when unexpected events arise.} \\
\\ 
\texttt{In addition to creating a schedule for parenting responsibilities, it's also essential to consider your child's extracurricular activities and schedule. This may involve coordinating with your co{-}parent to ensure that your child can participate in activities that are important to them. By creating a schedule that takes into account your child's interests and needs, you can help them develop a sense of stability and routine.} \\
\\ 
\texttt{Ultimately, creating a parenting plan that works for everyone involved requires careful consideration and planning. By taking the time to understand your daily schedule and creating a realistic plan, you can ensure that your child receives the care and attention they need, even when you're not together. Remember to be flexible and adapt to changing circumstances, and always prioritize your child's needs and schedule.} \\
            }
        }
    }
    & & \\ \\

    \theutterance \stepcounter{utterance}  
    & & & \multicolumn{2}{p{0.3\linewidth}}{
        \cellcolor[rgb]{0.9,0.9,0.9}{
            \makecell[{{p{\linewidth}}}]{
                \texttt{\tiny{[GM$|$GM]}}
                {'keywords:forbidden_words': {'forbidden_words': ['sleep', 'cook', 'feed']}, 'keywords:frequency': {'relation': 'at least', 'keyword': 'schedule', 'frequency': 6}}
            }
        }
    }
    & & \\ \\

    \theutterance \stepcounter{utterance}  
    & & & \multicolumn{2}{p{0.3\linewidth}}{
        \cellcolor[rgb]{0.9,0.9,0.9}{
            \makecell[{{p{\linewidth}}}]{
                \texttt{\tiny{[GM$|$GM]}}
                \texttt{game\_result = WIN} \\
            }
        }
    }
    & & \\ \\

\end{supertabular}
}

\end{document}
